\documentclass[12pt]{article}
\usepackage[utf8]{inputenc}
\usepackage[T1]{fontenc}
\usepackage[margin=1in]{geometry}
\usepackage{setspace}
\usepackage{xcolor}
\usepackage[colorlinks=true, urlcolor=blue!60!black, linkcolor=blue!60!black]{hyperref}
\setlength{\parindent}{2em}
\setlength{\parskip}{1.25em}
\renewcommand{\baselinestretch}{1.0}
\begin{document}
\begin{center}
{\LARGE \textit{Haaland v. Brackeen}, 599 U.S. 255 (2023)\par}
\vspace{0.8em}
{\large Case Brief\par}
\end{center}
\vspace{1em}

\section{Memory Jogger}
Rejecting a multi-front attack on the Indian Child Welfare Act, the Court held 7--2 that ICWA falls within Congress's Article~I power over Indian affairs, that its custody rules do not commandeer the States, and that the individual challengers and Texas lacked standing to press equal protection claims against federal officials (\href{https://supreme.justia.com/cases/federal/us/599/21-376/}{Justia}).

\section{Detailed Case Facts}
\subsection{The Indian Child Welfare Act of 1978}
Congress enacted the Indian Child Welfare Act (ICWA), 92 Stat.~3069, 25 U.S.C.~\S~1901 \textit{et seq.}, in 1978 after finding that ``an alarmingly high percentage of Indian families are broken up by the removal, often unwarranted, of their children from them by nontribal public and private agencies'' and that ``there is no resource that is more vital to the continued existence and integrity of Indian tribes than their children'' (\S~1901(3)--(4)) (\href{https://www.law.cornell.edu/supremecourt/text/21-376}{Cornell LII}). An ``Indian child'' is an unmarried minor who is either a tribal member or eligible for tribal membership and the biological child of a member (\S~1903(4)). Tribal courts hold exclusive jurisdiction over on-reservation Indian children (\S~1911(a)), while state courts exercise concurrent jurisdiction (subject to transfer) over off-reservation children (\S~1911(b)) and must apply ICWA's federal rules:
\begin{itemize}
    \item \textbf{Involuntary Proceedings (\S~1912):} ``Any party'' seeking involuntary foster care placement or termination of parental rights must notify the parent/custodian and tribe (\S~1912(a)), prove ``active efforts'' to prevent the breakup of the Indian family (\S~1912(d)), and establish via qualified expert testimony that continued parental custody is likely to cause ``serious emotional or physical damage''---by clear and convincing evidence for foster care (\S~1912(e)) and beyond a reasonable doubt for termination (\S~1912(f)).
    \item \textbf{Voluntary Proceedings and Intervention (\S\S~1911(c), 1913, 1914):} A tribe may intervene at any point and collaterally attack decrees violating ICWA, even against biological parents' wishes.
    \item \textbf{Placement Preferences (\S~1915(a)--(c)):} Absent ``good cause,'' adoptive placements prioritize (1) extended family, (2) other tribal members, or (3) ``other Indian families'' (\S~1915(a)); foster care follows a parallel four-tier hierarchy (\S~1915(b)). Under \S~1915(c), the child's tribe may alter the preference order by resolution.
    \item \textbf{Recordkeeping (\S\S~1915(e), 1951(a)):} States must maintain records of placement-preference compliance for inspection by the Secretary of the Interior or the tribe (\S~1915(e)), and state courts must transmit final adoption decrees and specified information to the Secretary (\S~1951(a)).
\end{itemize}

\subsection{Underlying Custody Disputes and Lower Court Rulings}
The suit arose from three custody proceedings (\href{https://www.law.cornell.edu/supremecourt/text/21-376}{Cornell LII}): (1) \textbf{The Brackeens (Texas)} fostered A.L.M. (Navajo/Cherokee) and sought to adopt him with his biological family's support; after the Navajo Nation proposed an unrelated Navajo couple in New Mexico and the state court denied adoption, the Brackeens sued, the Navajo couple withdrew, and the adoption closed, though they later sought to adopt A.L.M.'s sister, Y.R.J., over tribal objection; (2) \textbf{Altagracia Hernandez and the Librettis (Nevada/Texas)} sought to finalize Hernandez's voluntary placement of Baby O. with the Librettis after the Ysleta del Sur Pueblo Tribe enrolled Baby O. and intervened before withdrawing; and (3) \textbf{The Cliffords (Minnesota)} fostered Child P. and were denied adoption after the White Earth Band of Ojibwe reversed its eligibility determination, enrolled Child P., and Minnesota placed her with her maternal grandmother.

\section{Procedural History}
The individual petitioners joined Texas, Indiana, and Louisiana in suing federal officials in the Northern District of Texas; four tribes (Cherokee, Oneida, Quinault, and Morongo, later joined by the Navajo Nation) intervened as defendants (\href{https://supreme.justia.com/cases/federal/us/599/21-376/}{Justia}). The District Court granted plaintiffs summary judgment on equal protection, Tenth Amendment, nondelegation, and APA grounds. A Fifth Circuit panel reversed (937 F.3d 406 (2019)). Sitting en banc, the Fifth Circuit (994 F.3d 249 (2021) (per curiam)) upheld Congress's Article I power and rejected nondelegation, but struck down \S~1912(d), \S~1912(e)--(f) (experts), and \S~1915(e) under the Tenth Amendment, and divided evenly---affirming invalidation---on equal protection challenges to the ``other Indian families'' and ``Indian foster home'' preferences and Tenth Amendment challenges to \S~1912(a), \S~1915(a)--(b), and \S~1951(a). Only Texas remained among the States on certiorari.

\section{Judicial Votes}
The Supreme Court decided the case by a \textbf{7--2} vote. It was argued on November 9, 2022, and decided on June 15, 2023 (\href{https://supreme.justia.com/cases/federal/us/599/21-376/}{Justia}; \href{https://www.law.cornell.edu/supremecourt/text/21-376}{Cornell LII}). Disposition: 994 F.3d 249 affirmed in part (Article I legislative authority), reversed in part (Tenth Amendment anticommandeering), and vacated and remanded in part with instructions to dismiss for lack of jurisdiction (equal protection and nondelegation).

\begin{itemize}
    \item \textbf{Majority Opinion:} Justice Barrett delivered the opinion of the Court, joined by Chief Justice Roberts and Justices Sotomayor, Kagan, Gorsuch, Kavanaugh, and Jackson.
    \item \textbf{Concurring Opinions:} Justice Gorsuch filed a concurring opinion, joined by Justices Sotomayor and Jackson as to Parts I and III. Justice Kavanaugh filed a concurring opinion.
    \item \textbf{Dissenting Opinions:} Justice Thomas and Justice Alito each filed a dissenting opinion.
\end{itemize}

\section{Holding}
\textbf{Questions Presented:}

\begin{enumerate}
    \item \textbf{Article I Power:} Does Congress possess Article I authority over Indian affairs to enact ICWA's regulation of state child-custody proceedings notwithstanding traditional state control over family law?
    \item \textbf{Tenth Amendment Anticommandeering:} Do ICWA's involuntary-proceeding rules (\S~1912), placement preferences (\S~1915(a)--(b)), and recordkeeping mandates (\S\S~1915(e), 1951(a)) unconstitutionally commandeer state agencies or state courts?
    \item \textbf{Article III Standing:} Do the individual petitioners or Texas have standing in federal court against federal officials to bring Fifth Amendment equal protection or nondelegation (\S~1915(c)) challenges?
\end{enumerate}

\textbf{Holding:}

\begin{enumerate}
    \item \textbf{Yes (Article I Affirmed):} ICWA falls within Congress's broad Article I authority over Indian affairs, and the Constitution contains no categorical ``family law carveout'' (\href{https://www.law.cornell.edu/supremecourt/text/21-376}{Cornell LII}).
    \item \textbf{No (Tenth Amendment Reversed):} None of the challenged provisions commandeer the States. Section~1912 applies evenhandedly to private and state parties; \S~1915 imposes no search duty on state agencies and validly preempts state law in state courts; and \S\S~1915(e) and 1951(a) impose permissible ``ancillary'' judicial recordkeeping duties.
    \item \textbf{No (Standing Claims Vacated and Dismissed):} Neither the individual petitioners nor Texas has Article III standing to raise the equal protection or nondelegation claims.
\end{enumerate}

\section{Analysis of Opinions}
\subsection{Majority Opinion (Justice Barrett)}

\subsection{Article I Legislative Power (Part II)}
Justice Barrett surveyed precedent describing Congress's Indian-affairs authority as ``plenary and exclusive'' (\textit{United States v. Lara}, 541 U.S. 193, 200 (2004)), while stressing that ``plenary'' is not ``free-floating'' or unbounded. Precedent grounds this power in four sources: (1) the \textbf{Indian Commerce Clause} (Art.~I, \S~8, cl.~3), which reaches non-trade ``Indian affairs'' (\textit{Cotton Petroleum Corp. v. New Mexico}, 490 U.S. 163, 192 (1989)); (2) the \textbf{Treaty Clause} (Art.~II, \S~2, cl.~2); (3) \textbf{constitutional structure} and preconstitutional national powers (\textit{Morton v. Mancari}, 417 U.S. 535, 551--52 (1974)); and (4) the \textbf{federal--tribal trust relationship} (\textit{United States v. Mitchell}, 463 U.S. 206, 225--26 (1983)).

Although Congress lacks a general domestic-relations power, ``the Constitution does not erect a firewall around family law''; valid Article I legislation preempts conflicting state family law (\textit{Ridgway v. Ridgway}, 454 U.S. 46, 54 (1981); \textit{Fisher v. District Court}, 424 U.S. 382, 390 (1976)). Petitioners' narrower objections failed because they treated the slate as clean without asking to overrule precedent: \textit{United States v. Holliday}, 3 Wall.~407, 416--17 (1866), forecloses limiting the Indian Commerce Clause to tribes rather than individuals; \textit{Cotton Petroleum} answers the refrain that ``children are not commodities''; and \textit{Lara} refutes limiting structural powers to war and peace.

\subsection{Tenth Amendment Anticommandeering (Part III)}
\begin{enumerate}
    \item \textbf{Involuntary Proceedings (\S~1912):} Legislation applying ``evenhandedly'' to state and private actors ordinarily does not violate the Tenth Amendment (\textit{Murphy v. NCAA}, 584 U.S. 453 (2018); \textit{Reno v. Condon}, 528 U.S. 141 (2000)). Because \S~1912(d)'s ``active efforts'' rule (and notice, expert, and evidentiary rules) applies to ``any party''---including private litigants who regularly initiate involuntary proceedings---and \S~1922 exempts emergency removals, \S~1912 does not commandeer sovereign power (\href{https://www.law.cornell.edu/supremecourt/text/21-376}{Cornell LII}).
    \item \textbf{Placement Preferences (\S~1915):} Under \textit{Adoptive Couple v. Baby Girl}, 570 U.S. 637, 654 (2013), \S~1915 imposes no ``diligent search'' duty on state agencies; the objecting party or tribe must produce an alternative placement. Requiring \emph{state courts} to apply federal rules of decision---even in state-law causes of action---is valid Supremacy Clause preemption (\textit{New York v. United States}, 505 U.S. 144, 178--79 (1992)).
    \item \textbf{Recordkeeping (\S\S~1915(e), 1951(a)):} Confirming \textit{Printz v. United States}, 521 U.S. 898, 907--08 \& n.~2 (1997), early federal statutes (including the 1790--1802 Naturalization Acts, the Judiciary Act of 1789, and the 1790 Seamen's Act) and the Madisonian Compromise establish that Congress may impose ``ancillary'' recordkeeping and reporting duties on state courts.
\end{enumerate}

\subsection{Article III Standing (Part IV)}
The \textbf{individual petitioners} alleged an ``unequal footing'' injury but lacked \emph{redressability}: state courts and agencies, not the federal defendants, enforce ICWA's placement preferences, so neither an injunction nor a declaratory judgment against federal officials binds nonparty state officials (\textit{Lujan v. Defenders of Wildlife}, 504 U.S. 555, 569 (1992)). Hoping state courts will defer to a federal opinion cannot satisfy Article III. \textbf{Texas} lacked standing because a State has no equal protection rights of its own (\textit{South Carolina v. Katzenbach}, 383 U.S. 301, 323 (1966)), cannot sue the United States as \textit{parens patriae} (\textit{Alfred L. Snapp \& Son, Inc. v. Puerto Rico}, 458 U.S. 592, 610 n.~16 (1982)), cannot rely on an ``unclean hands'' or Spending Clause theory, and cannot trace its administrative costs to the independent placement preferences (\textit{California v. Texas}, 593 U.S. 659 (2021)). Both parties likewise lacked standing to challenge \S~1915(c) under nondelegation.

\subsection{Concurring Opinion (Justice Gorsuch)}

Justice Gorsuch chronicled the history of federal Indian boarding schools and mid-twentieth-century state child removals (Part I) and, writing alone in Part II, advanced an originalist ``Indian-law bargain'': Tribes retain inherent sovereignty, States have no authority over Indian affairs, and ``Commerce with the Indian Tribes'' was a Founding-era term of art covering bilateral diplomatic and noneconomic intercourse with individual Indians. Rejecting nineteenth-century ``plenary power'' rhetoric (\textit{United States v. Kagama}, 118 U.S. 375 (1886)) as an atextual misstep, he concluded ICWA validly shields tribal political survival from state intrusion (\href{https://www.law.cornell.edu/supremecourt/text/21-376}{Cornell LII}).

\subsection{Concurring Opinion (Justice Kavanaugh)}

Justice Kavanaugh emphasized that the equal protection challenge remains undecided and ``is serious,'' citing \textit{Palmore v. Sidoti}, 466 U.S. 429 (1984), and noting that courts can resolve it when properly raised by a litigant with standing in state-court custody proceedings (\href{https://www.law.cornell.edu/supremecourt/text/21-376}{Cornell LII}).

\subsection{Dissenting Opinion (Justice Thomas)}

Justice Thomas argued that the majority inverted enumerated-powers analysis, that the Founders rejected an ``Indian affairs'' clause, that ``Commerce'' carries the same economic meaning across all three Commerce Clauses, and that ``plenary power'' is unmoored ``judicial ipse dixit'' incapable of displacing state family law for off-reservation U.S. citizens (\href{https://www.law.cornell.edu/supremecourt/text/21-376}{Cornell LII}).

\subsection{Dissenting Opinion (Justice Alito)}

Justice Alito criticized the majority's claim that Congress's power is ``plenary'' yet not ``absolute'' as violating basic logic and contended that federal Indian authority cannot override foundational federalism limits or subordinate vulnerable children's ``best interests'' in state family courts (\href{https://www.supremecourt.gov/opinions/22pdf/21-376_7l48.pdf}{SupremeCourt.gov}).

\subsection{Subsequent Developments}

\begin{itemize}

\item \textbf{Redressability and State Standing:} Days later, \textit{United States v. Texas}, 599 U.S. 670 (2023) (\href{https://supreme.justia.com/cases/federal/us/599/22-58/}{Justia}), quoted \textit{Brackeen}'s rule that ``it is the judgment, not the opinion, that demonstrates redressability.'' In \textit{Murthy v. Missouri}, 603 U.S. 43 (2024) (\href{https://www.law.cornell.edu/supremecourt/text/23-411}{Cornell LII}), Justice Barrett again cited \textit{Brackeen}, 599 U.S. at 295, to confirm that States lack \textit{parens patriae} standing to sue the Federal Government on behalf of citizens. Meanwhile, Justice Thomas's concurrence in \textit{Arizona v. Navajo Nation}, 599 U.S. 555 (2023) (\href{https://www.law.cornell.edu/supremecourt/text/21-1484}{Cornell LII}), cited \textit{Brackeen} to caution against treating the general ``trust relationship'' as creating enforceable federal duties.

\end{itemize}

\section{Examples: Future Applications}
\begin{enumerate}

\item \textbf{Same-Side (Statute Upheld): A Federal Rule for Tribal Members' Estates.} Congress enacts a statute requiring ``any party'' petitioning a state probate court over an off-reservation tribal member's intestate estate to show ``active efforts'' to notify the tribe, apply a federal heirloom preference favoring tribal relatives, and have state probate clerks transmit inventories to the Interior Department. Under \textit{Brackeen}, there is no ``probate carveout'' to Article I Indian-affairs power; the ``any party'' rule is evenhanded; the preference validly preempts state law in state courts; clerk reporting is an ancillary judicial duty; and neither a State nor a surviving spouse suing federal officials in federal court has standing to raise an equal protection challenge.

\item \textbf{Same-Side (No Standing): An Adoption Agency Sues the Secretary.} A private adoption agency that places children with non-Indian foster families sues the Secretary of the Interior in federal district court, alleging that ICWA's foster-care preferences in \S~1915(b) disadvantage its families on the basis of race, and seeks a declaratory judgment and an injunction against the Secretary. Under \textit{Brackeen}, the suit must be dismissed for lack of redressability: state courts and state agencies, not the Secretary, apply the placement preferences, so a judgment against the Secretary would not bind the state judges who decide each placement. That a favorable federal opinion might persuade state courts does not count, because it is the judgment, not the opinion, that redresses an injury. The agency's families must raise the claim as parties in a state-court custody proceeding, as Justice Kavanaugh's concurrence noted.

\item \textbf{Opposite-Side (Statute Invalid): A Nationwide Child-Welfare Mandate Under the Interstate Commerce Clause.} Now suppose Congress invokes the \emph{Interstate} Commerce Clause to enact a nationwide child-welfare law applying to \emph{all} children that directs ``every State's Department of Child Protective Services'' (and no private actors) to conduct 50-state extended-family searches and orders state legislatures to adopt federal foster-licensing codes. Here, the outcome flips on every axis: (1) under the Interstate Commerce Clause (\textit{Lopez}; \textit{Morrison}), non-economic family custody proceedings lack a commercial nexus, making ``children are not commodities'' dispositive; (2) singling out state agencies and state legislatures directly violates \textit{Printz}, \textit{New York}, and \textit{Murphy}; and (3) the State possesses direct sovereign standing to challenge the federal command.

\item \textbf{Opposite-Side (Provision Likely Invalid): A Search Duty Imposed Only on State Agencies.} Amending ICWA itself, Congress requires each State's child-welfare agency, and no other party, to complete a documented nationwide search for extended-family and tribal placements before any Indian child may enter foster care, and to file monthly search reports with the Secretary of the Interior. Congress's Article~I power over Indian affairs still covers the subject, so the challenge turns on the Tenth Amendment. \textit{Brackeen} upheld \S~1912 because it binds ``any party'' evenhandedly and upheld \S~1915 because, under \textit{Adoptive Couple v. Baby Girl}, it imposes no search duty on state agencies. A duty addressed only to state executive agencies, with reporting by agency officials rather than ancillary court recordkeeping, is the kind of direct command to state officers that \textit{Printz} and \textit{Murphy} forbid, so the State would likely prevail.

\item \textbf{Fence-Sitter (Undecided): An Equal Protection Challenge Raised in State Court.} Non-Indian foster parents seek to adopt a child who is eligible for membership in one tribe; a family belonging to a different tribe asks for the placement under \S~1915(a)'s ``other Indian families'' preference. In the state adoption proceeding, where the preference is actually applied, the foster parents argue that it is a racial classification barred by \textit{Palmore v. Sidoti}, 466 U.S. 429 (1984). Their standing is secure, so a court must reach the merits \textit{Brackeen} left open. The tribe answers that, under \textit{Morton v. Mancari}, 417 U.S. 535 (1974), classifications tied to tribal membership are political, not racial. The foster parents reply that a preference for families of any tribe looks less like a political tie to the child's own tribe and more like ancestry. Justice Kavanaugh called the issue ``serious,'' and the Court has not resolved it.

\end{enumerate}

\section{Critique}
\textbf{Unstated Outer Limits:} While conceding that Indian-affairs precedent is ``unwieldy,'' the majority declined to define the borders of Congress's ``sizeable sphere,'' relying on petitioners' failure to ask to overrule precedent---a stark contrast to \textit{United States v. Lopez}, 514 U.S. 549 (1995).

\textbf{Formalism in Evenhandedness:} Critics warn that allowing Congress to bypass \textit{Printz} simply by addressing ``any party'' elevates statutory phrasing over reality where state agencies bring most child-welfare proceedings.

\textbf{Competing Originalisms:} Justices Gorsuch and Thomas both marshaled Founding-era dictionaries and early statutes yet reached opposite conclusions on whether ``Commerce with the Indian Tribes'' was a specialized diplomatic term of art or a uniform economic term.

\section{Key Quotations}
\begin{quote}
``To be clear, however, `plenary' does not mean `free-floating.' A power unmoored from the Constitution would lack both justification and limits. So like the rest of its legislative powers, Congress's authority to regulate Indians must derive from the Constitution, not the atmosphere.''
\end{quote}
\noindent --- Barrett, J. (\href{https://www.law.cornell.edu/supremecourt/text/21-376}{Cornell LII}).

\begin{quote}
``It could not be otherwise---Article I gives Congress a series of enumerated powers, not a series of blank checks. Thus, we reiterate that Congress's authority to legislate with respect to Indians is not unbounded. It is plenary within its sphere, but even a sizeable sphere has borders. \ldots{} But the Constitution does not erect a firewall around family law.''
\end{quote}
\noindent --- Barrett, J. (\href{https://www.law.cornell.edu/supremecourt/text/21-376}{Cornell LII}).

\begin{quote}
``It is a federal court's judgment, not its opinion, that remedies an injury; thus it is the judgment, not the opinion, that demonstrates redressability. The individual petitioners can hope for nothing more than an opinion, so they cannot satisfy Article III.''
\end{quote}
\noindent --- Barrett, J. (\href{https://www.law.cornell.edu/supremecourt/text/21-376}{Cornell LII}).

\begin{quote}
``Our Constitution reserves for the Tribes a place---an enduring place---in the structure of American life. It promises them sovereignty for as long as they wish to keep it. And it secures that promise by divesting States of authority over Indian affairs and by giving the federal government certain significant (but limited and enumerated) powers aimed at building a lasting peace.''
\end{quote}
\noindent --- Gorsuch, J., concurring (\href{https://www.law.cornell.edu/supremecourt/text/21-376}{Cornell LII}).


\end{document}
